\chapter{Lung biopsy}
\index{Lung biopsy}

\section*{Definition and Overview}
A lung biopsy is a diagnostic medical procedure in which a sample of lung tissue or cells is obtained for histopathological, cytological, or microbiological examination. It is performed when less invasive diagnostic modalities, such as sputum cytology or imaging, are insufficient to characterize a pulmonary abnormality. A lung biopsy is crucial for differentiating malignant lung nodules from benign infectious or inflammatory conditions, and for diagnosing diffuse parenchymal lung diseases. The choice of biopsy technique depends on the location and size of the lesion, the suspected pathology, and the patient's functional status. Modalities include percutaneous core needle biopsy (typically guided by computed tomography [CT] or ultrasound), flexible bronchoscopy (often combined with endobronchial ultrasound [EBUS]), and surgical lung biopsy (typically via video-assisted thoracoscopic surgery [VATS] or, rarely, open thoracotomy).

\section*{Epidemiology}
Lung biopsies are frequently performed worldwide, primarily driven by the high prevalence of pulmonary nodules detected incidentally on thoracic imaging or during structured lung cancer screening programmes. Lung cancer is the leading cause of oncology-related mortality globally and in many countries, necessitating accurate tissue diagnosis and molecular staging. The demographics of patients undergoing lung biopsy correspond closely with those of lung cancer and chronic lung diseases, typically involving older adults (aged 50 and above) with a history of tobacco smoking or occupational exposure to carcinogens (such as asbestos or silica). Percutaneous needle biopsy has a pneumothorax complication rate of approximately 15--20\%, with 5--7\% of cases requiring a chest tube, making it a common but risk-associated outpatient procedure.

\section*{Aetiology and Pathophysiology}
The primary clinical indications for performing a lung biopsy, along with the underlying diagnostic physiological goals, include:
\begin{itemize}
    \item \textbf{Investigation of pulmonary nodules:} Determining the nature of solitary or multiple lung nodules or masses, distinguishing malignant primary or metastatic lesions from benign granulomas or hamartomas.
    \item \textbf{Diagnosis of interstitial lung disease (ILD):} Identifying specific histopathological patterns (e.g. usual interstitial pneumonia [UIP], non-specific interstitial pneumonia [NSIP], or sarcoidosis) when high-resolution chest CT findings are non-diagnostic.
    \item \textbf{Identification of atypical infections:} Confirming opportunistic infections (such as fungal, mycobacterial, or \textit{Pneumocystis jirovecii} pneumonias) in immunocompromised patients when sputum samples are non-diagnostic.
    \item \textbf{Lung cancer staging and molecular profiling:} Sampling mediastinal lymph nodes via EBUS-guided transbronchial needle aspiration (EBUS-TBNA) to guide staging, and obtaining adequate tissue for biomarker testing (e.g. EGFR, ALK, ROS1, and PD-L1 expression) to direct targeted therapies.
    \item \textbf{Evaluation of treatment resistance:} Determining whether a previously diagnosed malignancy has mutated or transformed, or investigating persistent, unresolved pulmonary infiltrates.
\end{itemize}

\section*{Clinical Features}
While the biopsy itself is a transient procedure, the patient presents with symptoms of the underlying lung disease and experiences specific, expected post-procedural features:
\begin{itemize}
    \item \textbf{Pre-procedural symptoms:} Patients may be entirely asymptomatic (incidental radiological findings) or present with respiratory symptoms including a chronic cough, dyspnoea, haemoptysis, or pleuritic chest pain.
    \item \textbf{Post-procedural chest wall pain:} Localised, dull, or aching chest pain at the needle insertion site (percutaneous biopsy) or surgical incision sites (VATS), which is usually transient and responsive to simple analgesics.
    \item \textbf{Post-procedural cough and mild haemoptysis:} A mild, self-limiting cough, occasionally producing blood-tinged sputum or minor streaks of blood, is common for 24--48 hours after bronchoscopic or percutaneous biopsy.
    \item \textbf{Pleuritic chest pain:} Sharp, stabbing pain that worsens on deep inspiration, which can indicate pleural irritation or the development of a pneumothorax.
    \item \textbf{Mild respiratory distress:} Transient shortness of breath immediately following the procedure, requiring oxygen supplementation and monitoring of peripheral oxygen saturation.
\end{itemize}

\section*{Differential Diagnosis}
In the post-procedural monitoring period, clinicians must distinguish expected post-biopsy symptoms from acute, potentially life-threatening complications:
\begin{itemize}
    \item \textbf{Pneumothorax vs. musculoskeletal pain:} Distinguishing the sudden, sharp, unilateral pleuritic pain and progressive dyspnoea of a pneumothorax from normal chest wall soreness at the biopsy site.
    \item \textbf{Severe pulmonary haemorrhage vs. self-limiting haemoptysis:} Differentiating minor blood-tinged sputum from major endobronchial bleeding, which presents as continuous coughing of bright red blood, respiratory distress, and haemodynamic instability.
    \item \textbf{Post-procedural pneumonia vs. normal inflammatory response:} Distinguishing persistent high fever, purulent sputum production, and worsening cough from a mild, transient post-procedure inflammatory pyrexia.
    \item \textbf{Pulmonary embolism vs. procedure-induced pleuritis:} Differentiating acute dyspnoea, tachycardia, and pleuritic pain due to a thromboembolic event from localized pleural irritation caused by the biopsy needle or surgical scope.
\end{itemize}

\section*{When to Seek Medical Care}
Patients must be educated on post-procedure warning signs. They should contact their treating team or seek urgent medical attention if they develop worsening shortness of breath, high fevers, or persistent chest pain.

\textbf{Contact your local emergency services for:}
\begin{itemize}
    \item Sudden, severe, or worsening shortness of breath or difficulty breathing.
    \item Sharp, stabbing chest pain that makes it difficult to take a breath.
    \item Coughing up significant amounts of bright red blood.
    \item Blueness of the lips, face, or fingernails (cyanosis), or sudden collapse.
\end{itemize}

\section*{Diagnosis and Investigations}
The diagnosis is planned and confirmed through extensive pre-procedural imaging and post-procedural histopathological analysis:
\begin{itemize}
    \item \textbf{High-resolution chest CT scan:} The primary imaging modality used to locate the target lesion, plan the safest entry trajectory, and determine the optimal biopsy method.
    \item \textbf{Coagulation screen and platelet count:} Mandatory pre-procedural laboratory tests to rule out bleeding diatheses (aiming for INR $<1.5$ and platelets $>50 \times 10^9$/L) to minimise the risk of severe intra-pulmonary haemorrhage.
    \item \textbf{Pulmonary function tests (PFTs):} Evaluates FEV1 and DLCO to assess the patient's respiratory reserve and their ability to tolerate a potential pneumothorax or lung resection.
    \item \textbf{Post-procedural chest X-ray:} A routine erect chest X-ray is performed 1 to 4 hours following percutaneous needle or transbronchial biopsy to screen for the presence of a pneumothorax or clinical haemorrhage.
    \item \textbf{Histopathological and cytological analysis:} Microscopic evaluation of core tissue or cell blocks using special stains (e.g. H\&E, Ziehl-Neelsen for mycobacteria, Grocott's methenamine silver for fungi) and immunohistochemical panels to establish a definitive diagnosis.
\end{itemize}

\section*{Management}
The management of a patient undergoing a lung biopsy involves the execution of the chosen technique, post-procedure observation, and the management of any immediate complications:
\begin{itemize}
    \item \textbf{Procedure execution:} Utilising local anaesthesia and conscious sedation for percutaneous biopsies, or general anaesthesia for bronchoscopic EBUS and surgical VATS procedures.
    \item \textbf{Post-procedure monitoring:} Assessing vital signs, respiratory effort, and oxygen saturation every 15--30 minutes in a recovery unit for a minimum of 2--4 hours.
    \item \textbf{Pneumothorax management:} Asymptomatic or small pneumothoraces ( $<2$ cm) are managed conservatively with bed rest, high-flow supplemental oxygen (which accelerates pleural air absorption), and serial chest X-rays. Large or symptomatic pneumothoraces require urgent insertion of a chest tube (intercostal catheter) or a small-bore pleural aspiration catheter.
    \item \textbf{Haemostasis:} For mild endobronchial bleeding during bronchoscopy, instillation of cold saline, topical adrenaline, or balloon tamponade may be employed.
    \item \textbf{Analgesia:} Administration of paracetamol or mild opioids to manage post-procedural pain, avoiding non-steroidal anti-inflammatory drugs (NSAIDs) in the immediate post-operative window due to bleeding risk.
\end{itemize}

\section*{Complications}
The potential complications of a lung biopsy vary by the technique used but generally include:
\begin{itemize}
    \item \textbf{Pneumothorax:} Air leaking into the pleural cavity, which is the most common complication of percutaneous needle biopsy.
    \item \textbf{Pulmonary haemorrhage:} Bleeding into the lung tissue or airways, which can range from minor localized bruising to severe, life-threatening alveolar haemorrhage.
    \item \textbf{Haemothorax:} Accumulation of blood within the pleural space due to laceration of an intercostal or pulmonary blood vessel.
    \item \textbf{Air embolism:} A rare but catastrophic event where air enters the pulmonary circulation and travels to the heart or brain, causing stroke or myocardial infarction.
    \item \textbf{Infection:} Localised wound infection, pneumonia, or empyema (pus in the pleural cavity).
    \item \textbf{Tumour seeding:} Extremely rare implantation of tumor cells along the needle tract during a percutaneous biopsy.
\end{itemize}

\section*{Prognosis and Outcomes}
The prognosis and diagnostic utility of a lung biopsy are highly favourable. CT-guided percutaneous biopsy and EBUS-TBNA have diagnostic yields exceeding 90\% for malignant lesions. VATS surgical lung biopsy approaches 100\% diagnostic accuracy for both malignant and interstitial lung diseases. Most patients undergoing outpatient percutaneous or bronchoscopic biopsies are discharged home on the same day and recover fully within 24--48 hours. Patients who undergo surgical VATS biopsies require a hospital stay of 1--3 days (often with a temporary chest drain) and recover fully within 1--2 weeks. The long-term prognosis is determined entirely by the underlying pathology diagnosed (e.g. stage and type of lung cancer, or severity of interstitial lung disease).

\section*{Prevention and Self-Care}
To minimise risks and support healing, patients must adhere to specific post-procedure self-care instructions:
\begin{itemize}
    \item \textbf{Activity restriction:} Avoid strenuous exercise, heavy lifting (over 5 kg), or vigorous physical activity for at least 48 hours after a needle biopsy, and for up to 2 weeks after surgical VATS.
    \item \textbf{Travel restrictions:} Strictly avoid commercial air travel or scuba diving for at least 1 to 2 weeks post-procedure, or until a chest X-ray confirms that any pneumothorax has completely resolved.
    \item \textbf{Cough management:} Take prescribed cough suppressants to prevent sudden increases in intrathoracic pressure, which can provoke bleeding or a pneumothorax.
    \item \textbf{Wound hygiene:} Keep the biopsy site clean and dry. Gently wash the area with soap and water after 24 hours, and avoid swimming or soaking in baths until the wound is closed.
    \item \textbf{Smoking cessation:} Cease smoking immediately to improve pulmonary recovery, enhance wound healing, and reduce the risk of future respiratory complications.
\end{itemize}

\section*{Sources and Further Reading}
The following reputable organisations provide further information on lung biopsies:
\begin{itemize}
    \item Lung Foundation Australia. \textit{Patient resources on lung procedures and diagnostics.} https://lungfoundation.com.au/
    \item Thoracic Society of Australia and New Zealand (TSANZ). \textit{Clinical statements and guidance on pleural and lung interventions.} https://www.thoracic.org.au/
    \item NHS. \textit{Overview of lung biopsy, including needle and surgical procedures.} https://www.nhs.uk/
    \item Mayo Clinic. \textit{Information on CT-guided and bronchoscopic lung biopsy.} https://www.mayoclinic.org/
    \item American Thoracic Society (ATS). \textit{Patient education on bronchoscopy and lung procedures.} https://www.thoracic.org/
\end{itemize}
